\documentclass[a4paper,11pt]{article}
\usepackage[T1]{fontenc}
\usepackage[utf8]{inputenc}
\usepackage{geometry}\geometry{margin=2cm}
\usepackage{amsmath}
\usepackage{rootlocus}

\title{\textbf{rootlocus} --- Demonstration}
\date{2026}

\begin{document}
\maketitle

\section{$H(s) = \dfrac{1}{s^2+3s+2}$ --- 2 poles reels}

Poles : $s=-1$, $s=-2$. Pas de zeros.

\begin{center}
\rootlocus[width=8cm, height=6cm,
           label={$H(s)=\frac{1}{s^2+3s+2}$}]
          {1}{1,3,2}
\end{center}

\section{$H(s) = \dfrac{s+1}{s(s+2)(s+3)}$ --- 3 poles, 1 zero}

\begin{center}
\rootlocus[width=8cm, height=7cm,
           label={$H(s)=\frac{s+1}{s(s+2)(s+3)}$}]
          {1,1}{1,5,6,0}
\end{center}

\section{$H(s) = \dfrac{1}{s^2+2s+5}$ --- poles complexes conjugues}

Poles : $s = -1 \pm 2j$.

\begin{center}
\rootlocus[width=8cm, height=7cm,
           label={$H(s)=\frac{1}{s^2+2s+5}$}]
          {1}{1,2,5}
\end{center}

\end{document}
